\chapter{Representação do conhecimento}
\label{cap:representacao}

\section{Tipos de conhecimento}
\label{sec:tipos-conhecimento}

A base de conhecimento reúne as regras e relações necessárias para
interpretar a informação disponível sobre o estado atual do veículo e do
ambiente e determinar as condições que podem conduzir a diferentes propostas
de decisão. No projeto são considerados diferentes tipos de conhecimento.
Alguns correspondem diretamente à situação observada, como a velocidade
atual, a presença de um veículo à frente, a distância a um obstáculo, o
estado de um semáforo ou a ocupação de uma via de trânsito. Outros resultam
da interpretação ou combinação dessa informação, como o TTC ou a
classificação de uma situação segundo o respetivo nível de risco.

Importa distinguir a base de conhecimento da memória de trabalho. A primeira
contém as regras utilizadas pelo motor de inferência, enquanto a memória de
trabalho mantém os factos e estados que representam aquilo que o sistema
conhece num determinado momento. Neste sentido, a organização do
conhecimento considerado no projeto pode ser dividida em tipos, como
representado na \cref{tab:tipos-conhecimento}.

\begin{table}[H]
  \centering
  \caption{Tipo de conhecimento e exemplos utilizados}
  \label{tab:tipos-conhecimento}
  \small
  \begin{tabularx}{\textwidth}{@{}L{3.6cm}VX@{}}
    \toprule
    \textbf{Tipo de conhecimento} & \textbf{Exemplos} \\
    \midrule
    Entidades & veículo, peão, obstáculo, via de trânsito, sinal de trânsito \\ \linhasep
    Propriedades & velocidade, distância, meteorologia, velocidade relativa \\ \linhasep
    Estados & veículo em movimento, via ocupada, semáforo vermelho, chuva
      presente \\ \linhasep
    Dados derivados & TTC e outros indicadores calculados a partir das leituras
      disponíveis \\ \linhasep
    Factos semânticos & risco alto, situação crítica, distância insegura \\ \linhasep
    Ações & manter velocidade, acelerar, desacelerar, travar, travar em
      emergência, parar, manter a via, mudar de via \\ \linhasep
    Níveis de risco & baixo, moderado, alto e crítico \\ \linhasep
    Prioridades na decisão & progressão e conforto, cumprimento da velocidade,
      segurança na manobra, perigo rodoviário, prevenção da colisão, paragem
      legal obrigatória e segurança e emergência \\
    \bottomrule
  \end{tabularx}
\end{table}

Para além desta organização, há diferença na natureza da informação
utilizada, nomeadamente:
\begin{itemize}
  \item conhecimento factual, associado ao estado atual do veículo e do
    ambiente;
  \item conhecimento temporal, relacionado com os eventos e a evolução de
    determinadas observações;
  \item conhecimento normativo, associado às regras como, por exemplo, a
    obrigação de parar tendo em conta determinados sinais;
  \item conhecimento heurístico, utilizado em situações de avaliação de risco;
  \item metaconhecimento, com o objetivo de definir uma ordem nas propostas de
    decisão.
\end{itemize}

A informação associada à razão apresentada, à regra de origem e à evidência
utilizada na tomada de decisão é, igualmente, conhecimento a considerar
aquando da explicação do processo de inferência.

\section{Estrutura das regras}
\label{sec:estrutura-regras}

Esta distinção permite separar aquilo que é diretamente observado daquilo que
é concluído pelo sistema. Por exemplo, a distância a um veículo da frente e a
respetiva velocidade relativa constituem informação disponível para a
inferência, enquanto a conclusão de que existe uma situação de risco elevado
resulta da análise desses elementos através das regras definidas na base de
conhecimento. Depois de representada sob a forma de factos, estados e
eventos, esta informação é utilizada pelo motor de inferência. De forma
simplificada, uma regra pode ser representada através da seguinte estrutura:

\begin{tcolorbox}[enhanced, colback=corFundo, colframe=corEscura,
    boxrule=0pt, leftrule=1.2pt, arc=0pt, outer arc=0pt,
    left=8pt, right=8pt, top=4pt, bottom=4pt,
    width=0.6\textwidth, center, fontupper=\ttfamily\small]
  SE condição 1\\
  E condição 2\\
  E condição 3\\
  ENTÃO ação proposta\\
  NÍVEL DE RISCO = nível\\
  PRIORIDADE DA DECISÃO = categoria
\end{tcolorbox}

Neste sentido, uma regra pode depender de vários factos em simultâneo. A
título de exemplo, detetar um peão não conduz, necessariamente, à mesma ação
em todas as situações, uma vez que também podem ser relevantes a distância, a
velocidade do veículo e o nível de risco identificado. Outro exemplo é a
mudança de via, que poderá ser impedida caso a via de trânsito adjacente
esteja ocupada ou exista uma decisão de segurança com prioridade superior.
Assim, uma decisão pode resultar da combinação de diferentes factos, dados
calculados e regras aplicáveis a uma situação concreta. Esta forma de
representação permite manter o conhecimento do domínio explícito e separado
da restante implementação.

Ficam fora do âmbito da aquisição de conhecimento e da respetiva
representação aspetos como o planeamento do percurso do veículo, o uso de
dados de sensores reais, o controlo detalhado da dinâmica do veículo, as
regras que permitam a cobertura de todas as situações rodoviárias possíveis e
a certificação necessária para a utilização num veículo real, uma vez que o
foco do trabalho está na aquisição, representação, inferência e explicação do
conhecimento.

\section{Regras representativas}
\label{sec:regras-representativas}

A \cref{tab:regras} apresenta algumas das regras mais representativas da
base de conhecimento, descritas em linguagem natural a partir das regras
implementadas no protótipo e das indicações do perito ainda por implementar.
Para cada regra, indicam-se as condições, a ação proposta, o nível de risco,
o valor da categoria de prioridade (\cref{sec:prioridades}), o tipo de
conhecimento que representa e o seu estado. Os limiares indicados são valores
provisórios definidos pela equipa, que deverão ser substituídos pelos
critérios discutidos com o perito (\cref{cap:sensores}).

\begin{table}[H]
  \centering
  \caption{Regras representativas da base de conhecimento}
  \label{tab:regras}
  \footnotesize
  \renewcommand{\arraystretch}{1.12}
  \begin{tabularx}{\textwidth}{@{}L{0.5cm}V
      >{\raggedright\arraybackslash}XV
      >{\raggedright\arraybackslash}p{3.1cm}V
      L{1.4cm}V C{0.95cm}V L{1.75cm}V L{1.95cm}@{}}
    \toprule
    \textbf{ID} & \textbf{SE} & \textbf{ENTÃO} & \textbf{Risco}
      & \textbf{Prior.} & \textbf{Tipo} & \textbf{Estado} \\
    \midrule
    R1 & velocidade acima do limite detetado
       & reduzir para o limite & moderado & 20 & normativo & implementada \\ \linhasep
    R2 & chuva
       & reduzir a velocidade, mais ainda se a visibilidade ou a aderência
         forem reduzidas & moderado & 20 & heurístico & implementada \\ \linhasep
    R3 & chuva e TTC inferior a 3~s em pelo menos 2 leituras nos últimos 2~s
       & travar até 15~km/h & alto & 50 & temporal & implementada \\ \linhasep
    R4 & peão na trajetória entre 6 e 24~m
       & parar & alto & 50 & heurístico & implementada \\ \linhasep
    R5 & peão na trajetória a 6~m ou menos
       & travagem de emergência & crítico & 70 & heurístico & implementada \\ \linhasep
    R6 & semáforo vermelho a 45~m ou menos
       & parar & alto & 60 & normativo & implementada \\ \linhasep
    R7 & pedido de mudança de via e via de destino ocupada
       & manter a via & moderado & 30 & normativo & implementada \\ \linhasep
    R8 & falha de um sensor (ausência de dados)
       & paragem controlada & crítico & 70 & heurístico & prevista \\ \linhasep
    R9 & distância ao objeto inferior à distância de travagem
         (\cref{eq:dtrav})
       & desvio, se a via adjacente estiver livre, ou travagem de emergência
       & crítico & 70 & heurístico & prevista \\
    \bottomrule
  \end{tabularx}
\end{table}

